\documentclass[a4,11pt]{aleph-notas}
% Se puede ver la documentación aquí:
% https://github.com/alephsub0/LaTeX_aleph-notas

% -- Paquetes adicionales
\usepackage{enumitem}
\usepackage{aleph-comandos}
\usepackage{booktabs}


% -- Datos
\institucion{Facultad de Ciencias Exactas, Naturales y Ambientales}
\carrera{Catálogo STEM}
\asignatura{Métodos Numéricos}
\tema{Resumen no. 7: Optimización numérica}
\autor{Mario Cueva - Andrés Merino}
\fecha{Periodo 2026-1}

\logouno[0.18\textwidth]{Logos/logoPUCE_04_ac}
\definecolor{colortext}{HTML}{1957d1}
\definecolor{colordef}{HTML}{1957d1}
\fuente{montserrat}


% -- Comandos adicionales
\setlist[enumerate]{label=\roman*.}
\newcommand{\vect}[1]{\mathbf{#1}}


\begin{document}

\encabezado

%%%%%%%%%%%%%%%%%%%%%%%%%%%%%%%%%%%%%%
\section{Optimización numérica: búsqueda exhaustiva y sección áurea}

\begin{defi}[Búsqueda exhaustiva]
    La \textbf{búsqueda exhaustiva} aproxima el mínimo de una función $f$ en un intervalo $[a,b]$ evaluando $f$ en $n+1$ puntos equidistantes
    \[
        x_i = a + i\,\frac{b-a}{n}, \qquad i=0,1,\ldots,n,
    \]
    y devolviendo el punto $x^*=x_k$ tal que $f(x_k)\leq f(x_i)$ para todo $i$.
\end{defi}

\begin{ejem}
    Para minimizar $f(x)=x^2-4x+5$ en $[0,4]$ con $n=4$ subintervalos, los nodos y valores son:
    \[
        \begin{array}{c|ccccc}
            x_i & 0 & 1 & 2 & 3 & 4\\
            \hline
            f(x_i) & 5 & 2 & 1 & 2 & 5
        \end{array}
    \]
    El menor valor es $f(2)=1$, por lo que la búsqueda devuelve $x^*=2$, que coincide con el mínimo exacto.
\end{ejem}

\begin{defi}[Método de la sección áurea]
    El \textbf{método de la sección áurea} reduce iterativamente el intervalo $[a,b]$ que encierra un mínimo. En cada iteración se evalúa $f$ en los dos puntos interiores
    \[
        x_1 = b-\varphi(b-a),
        \qquad
        x_2 = a+\varphi(b-a),
    \]
    donde $\varphi=\dfrac{\sqrt{5}-1}{2}\approx 0.618$ es la razón áurea. Luego:
    \begin{enumerate}
        \item Si $f(x_1)<f(x_2)$, el nuevo intervalo es $[a,\,x_2]$.
        \item Si $f(x_1)\geq f(x_2)$, el nuevo intervalo es $[x_1,\,b]$.
    \end{enumerate}
    En cada iteración el intervalo se reduce por el factor $\varphi$, por lo que el método converge a una tasa constante.
\end{defi}

\begin{ejem}
    Para minimizar $f(x)=x^2-2x+3$ en $[0,3]$ (mínimo exacto en $x=1$):
    \medskip

    \textbf{Iteración 1:} Con $a=0$, $b=3$ y $\varphi\approx 0.618$,
    \[
        x_1=3-0.618(3)=1.146,\qquad x_2=0+0.618(3)=1.854.
    \]
    \[
        f(1.146)=(1.146)^2-2(1.146)+3=2.021,\qquad f(1.854)=(1.854)^2-2(1.854)+3=2.729.
    \]
    Como $f(x_1)<f(x_2)$, el nuevo intervalo es $[0,\,1.854]$.
    \medskip

    \textbf{Iteración 2:} Con $a=0$, $b=1.854$,
    \[
        x_1=1.854-0.618(1.854)=0.708,\qquad x_2=0+0.618(1.854)=1.146.
    \]
    \[
        f(0.708)=(0.708)^2-2(0.708)+3=2.085,\qquad f(1.146)=2.021.
    \]
    Como $f(x_2)<f(x_1)$, el nuevo intervalo es $[0.708,\,1.854]$. El método continúa reduciendo el intervalo hacia $x=1$.
\end{ejem}

\begin{advertencia}
    La búsqueda exhaustiva requiere muchas evaluaciones de $f$ para lograr buena precisión. La sección áurea es más eficiente: con cada par de evaluaciones reduce el intervalo por el factor $\varphi\approx 0.618$, independientemente de la función.
\end{advertencia}


%%%%%%%%%%%%%%%%%%%%%%%%%%%%%%%%%%%%%%
\section{Optimización numérica: Método de Newton y Brent}

\begin{defi}[Método de Newton para optimización en varias variables]
    El \textbf{método de Newton para optimización} en $\mathbb{R}^n$ resuelve iterativamente $\nabla f({x})={0}$. A partir de una aproximación ${x}^{(k)}$, la actualización es
    \[
        {x}^{(k+1)}
        =
        {x}^{(k)}
        -
        \bigl[H_f\bigl({x}^{(k)}\bigr)\bigr]^{-1}
        \nabla f\bigl({x}^{(k)}\bigr),
    \]
    donde $H_f({x}^{(k)})$ es la matriz hessiana evaluada en ${x}^{(k)}$. Si el método converge a ${x}^*$ y $H_f({x}^*)$ es definida positiva, entonces ${x}^*$ es un mínimo local.
\end{defi}

\begin{ejem}
    Para minimizar $f(x,y)=x^2+2y^2-2xy$, se tienen
    \[
        \nabla f(x,y)=
        \begin{pmatrix}
            2x-2y\\
            4y-2x
        \end{pmatrix},
        \qquad
        H_f=
        \begin{pmatrix}
            2 & -2\\
            -2 & 4
        \end{pmatrix},
        \qquad
        H_f^{-1}=
        \begin{pmatrix}
            1 & \tfrac{1}{2}\\[1mm]
            \tfrac{1}{2} & \tfrac{1}{2}
        \end{pmatrix}.
    \]
    Partiendo de ${x}^{(0)}=(2,1)^T$, se calcula $\nabla f(2,1)=(2,0)^T$ y
    \[
        {x}^{(1)}
        =
        \begin{pmatrix}
            2\\1
        \end{pmatrix}
        -
        \begin{pmatrix}
            1 & \tfrac{1}{2}\\[1mm]
            \tfrac{1}{2} & \tfrac{1}{2}
        \end{pmatrix}
        \begin{pmatrix}
            2\\0
        \end{pmatrix}
        =
        \begin{pmatrix}
            2\\1
        \end{pmatrix}
        -
        \begin{pmatrix}
            2\\1
        \end{pmatrix}
        =
        \begin{pmatrix}
            0\\0
        \end{pmatrix}.
    \]
    El método converge en una sola iteración al mínimo global $(0,0)$ (resultado exacto porque $f$ es cuadrática).
\end{ejem}

\begin{defi}[Método de Brent]
    El \textbf{método de Brent} para minimización en una dimensión combina la robustez de la sección áurea con la rapidez de la \textbf{interpolación parabólica}. Dados tres puntos $x_1<x_2<x_3$ con $f(x_2)<f(x_1)$ y $f(x_2)<f(x_3)$, el mínimo de la parábola que los interpola es
    \[
        x^*
        =
        x_2
        -
        \frac{1}{2}
        \cdot
        \frac{(x_2-x_1)^2\bigl[f(x_2)-f(x_3)\bigr]-(x_2-x_3)^2\bigl[f(x_2)-f(x_1)\bigr]}{(x_2-x_1)\bigl[f(x_2)-f(x_3)\bigr]-(x_2-x_3)\bigl[f(x_2)-f(x_1)\bigr]}.
    \]
    Si $x^*$ cae dentro del intervalo y reduce suficientemente el dominio, se acepta; en caso contrario se usa un paso de sección áurea como respaldo.
\end{defi}

\begin{ejem}
    Para minimizar $f(x)=x^2$ en $[-1,2]$, se eligen los puntos iniciales $x_1=-1$, $x_2=0.5$, $x_3=2$, con valores $f(-1)=1$, $f(0.5)=0.25$, $f(2)=4$. Aplicando la fórmula:
    \[
        x^*
        =
        % 0.5
        % -
        % \frac{1}{2}
        % \cdot
        % \frac{(1.5)^2(-3.75)-(−1.5)^2(-0.75)}{(1.5)(-3.75)-(-1.5)(-0.75)}
        =
        0.5
        -
        \frac{1}{2}
        \cdot
        \frac{-6.75}{-6.75}
        =
        0.
    \]
    El método de Brent obtiene el mínimo exacto $x^*=0$ en una sola iteración gracias a la interpolación parabólica.
\end{ejem}

\begin{advertencia}
    El método de Brent garantiza convergencia (por la sección áurea como respaldo) y converge superlinealmente cerca del mínimo (por la interpolación parabólica). Es uno de los métodos más confiables para optimización en una sola variable.
\end{advertencia}


%%%%%%%%%%%%%%%%%%%%%%%%%%%%%%%%%%%%%%
\section{Optimización numérica con restricciones}

\begin{defi}[Multiplicadores de Lagrange]
    Para minimizar $f({x})$ sujeto a la restricción $g({x})=0$, el \textbf{método de multiplicadores de Lagrange} forma la función lagrangiana
    \[
        \mathcal{L}({x},\lambda)=f({x})-\lambda\,g({x}).
    \]
    Los puntos críticos del problema restringido satisfacen el sistema
    \[
        \nabla f({x}^*)=\lambda\,\nabla g({x}^*),
        \qquad
        g({x}^*)=0,
    \]
    donde el escalar $\lambda$ se denomina \textbf{multiplicador de Lagrange}.
\end{defi}

\begin{ejem}
    Para minimizar $f(x,y)=x^2+y^2$ sujeto a $g(x,y)=x+y-1=0$, la condición $\nabla f=\lambda\,\nabla g$ da
    \[
        (2x,\,2y)=\lambda(1,\,1),
    \]
    de donde $2x=\lambda$ y $2y=\lambda$, es decir, $x=y$. Sustituyendo en la restricción $x+y=1$ se obtiene $x=y=\tfrac{1}{2}$, con valor mínimo
    \[
        f\!\left(\tfrac{1}{2},\tfrac{1}{2}\right)=\frac{1}{4}+\frac{1}{4}=\frac{1}{2}.
    \]
\end{ejem}

\begin{defi}[Método de penalización]
    El \textbf{método de penalización} convierte un problema con restricciones en una sucesión de problemas sin restricciones. Para la restricción $g({x})=0$, se minimiza la función penalizada
    \[
        P_\mu({x})=f({x})+\mu\,\bigl[g({x})\bigr]^2,
    \]
    donde $\mu>0$ es el \textbf{parámetro de penalización}. Al hacer $\mu\to\infty$, la solución de $P_\mu$ converge a la del problema restringido original.
\end{defi}

\begin{ejem}
    Para el problema anterior, la función penalizada es
    \[
        P_\mu(x,y)=x^2+y^2+\mu(x+y-1)^2.
    \]
    Igualando las derivadas parciales a cero y explotando la simetría $x=y$ del sistema resultante, se obtiene
    \[
        x = y = \frac{\mu}{1+2\mu}.
    \]
    Para $\mu=10$ se tiene $x=y=\tfrac{10}{21}\approx 0.476$; para $\mu=100$ se tiene $x=y=\tfrac{100}{201}\approx 0.498$. Cuando $\mu\to\infty$, $x=y\to\tfrac{1}{2}$, recuperando la solución exacta.
\end{ejem}

\begin{advertencia}
    Al aumentar $\mu$ la función $P_\mu$ se vuelve más rígida y difícil de minimizar numéricamente. Una alternativa más robusta son los métodos de \textbf{Lagrangiano aumentado}, que combinan el término de penalización con actualizaciones explícitas del multiplicador $\lambda$.
\end{advertencia}


\end{document}
